% 08_busqueda_adversarial.tex - Capítulo 8: Búsqueda Adversarial (F4)
\chapter{Búsqueda Adversarial: Minimax y Alfa-Beta}
\label{ch:adversarial}

Este capítulo cubre la transición de ``encontrar una ruta'' a ``decidir frente a un adversario''
(RF-25 a RF-33). Se implementan Minimax y alfa-beta con profundidad configurable, utilidad
propia y 7 métricas RF-30.

\section{Cambio de Paradigma: De Búsqueda a Decisión}
\label{sec:cambio-paradigma}

En búsquedas F2/F3 (un agente): ``¿Cuál es mi mejor camino?'' $\to$ \code{PathProblem},
adversario ignorado.

En búsqueda adversarial F4: ``¿Cuál es mi mejor acción considerando las respuestas del
adversario?'' $\to$ \code{GameState} completo, \textbf{ambos bandos razonan}.

\section{Profundidad y Ply (G11, Decisión 2026-10-17)}
\label{sec:profundidad-ply}

\begin{itemize}
\item \textbf{1 unidad de profundidad = 1 acción de 1 bando} (1 ply)
\item Los bandos alternan por ply: MAX (lado que mueve en la raíz) $\to$ MIN $\to$ MAX $\to$ MIN...
\item \code{Wait} siempre es legal y cuenta como ply (consume el turno)
\item Profundidad configurable vía JSON: \code{juego.profundidad\_maxima\_minimax} o CLI \code{--depth}
\item Profundidad mínima = 1
\end{itemize}

\section{Función de Utilidad (RF-27, RF-33)}
\label{sec:utilidad}

Definida en \code{algorithms/utility.py}. Perspectiva: \textbf{MAX = bando que mueve en la raíz};
MAX maximiza, MIN minimiza. Pesos leídos del JSON (\code{utilidad}), defaults 10/3/1/0.5.

\begin{minted}{python}
U = w_entrega  * (+1 si MAX entregó, -1 si MIN entregó)
  + w_carga    * (+1 si MAX porta, -1 si MIN porta)
  - w_dist     * (distancia al objetivo actual)
  + w_rival    * (distancia del rival más cercano al portador, SOLO si MAX porta)
\end{minted}

\textbf{Objetivo dinámico (distancia):}
\begin{itemize}
\item MAX porta $\to$ objetivo = base propia de MAX
\item MIN porta $\to$ objetivo = posición del portador (perseguir al rival)
\item Nadie porta $\to$ objetivo = celda del recurso
\end{itemize}
Distancia = Manhattan desde la unidad MAX más cercana (desde el portador si MAX porta).

\textbf{Término rival:} Solo cuando MAX porta. Un rival cercano al portador de MAX \textbf{baja}
$U$ (riesgo de interceptación). Si MIN porta, no se suma (la persecución ya cubre ese riesgo).

\textbf{Pesos (desde JSON \code{utilidad}, defaults 10/3/1/0.5):}
\begin{table}[H]
\centering
\caption{Pesos de utilidad (RNF-03: configurables en JSON)}
\begin{tabular}{lll}
\toprule
\textbf{Peso} & \textbf{Default} & \textbf{Significado} \\
\midrule
\code{entrega} & 10 & Victoria/derrota domina todo \\
\code{recurso\_cargado} & 3 & Porta recurso > estar cerca \\
\code{distancia\_objetivo} & 1 & Cada paso de Manhattan = -1 \\
\code{distancia\_rival} & 0.5 & Rival cerca = riesgo (solo si MAX porta) \\
\bottomrule
\end{tabular}
\end{table}
\label{tab:pesos-utilidad}


\textbf{Terminal gana:} \code{winner()} devuelve $\pm$ \code{w\_entrega} ($\pm$10) y corta la
recursión. Victoria/derrota domina sobre cualquier combinación de distancias.

\section{Minimax y Alfa-Beta: Recursión Compartida}
\label{sec:minimax-alfabeta}

Ambos algoritmos comparten \textbf{una sola función recursiva} \code{value\_of};
alfa-beta solo añade el corte \code{alpha >= beta}.

\begin{minted}{python}
def value_of(node, level, alpha, beta):
    if level == depth or is_terminal(board, node):
        evaluated += 1
        return utility(node)
    maximizing = node.turn is max_side
    best = -inf if maximizing else inf
    for action in ordered(node):  # orden configurable: default/reversed
        child, _ = apply(board, node, action)
        generated += 1
        score = value_of(child, level + 1, alpha, beta)
        if maximizing:
            best = max(best, score)
            alpha = max(alpha, best)
        else:
            best = min(best, score)
            beta = min(beta, best)
        if prune and alpha >= beta:
            pruned += len(actions) - index - 1  # hermanos omitidos
            break
    return best
\end{minted}

\section{Métricas RF-30 (Implementadas en \code{AdversarialResult})}
\label{sec:metricas-adversarial}

\begin{table}[H]
\centering
\caption{Métricas adversariales (RF-30)}
\begin{tabular}{ll}
\toprule
\textbf{Métrica} & \textbf{Qué registra} \\
\midrule
\code{generated} & Hijos producidos (alfa-beta genera perezosamente: solo hasta corte) \\
\code{evaluated} & Estados donde se calculó utilidad (hojas + terminales) \\
\code{pruned} & Hermanos omitidos por \code{alpha >= beta} \\
\code{depth\_reached} & Nivel máximo alcanzado en la recursión \\
\code{time\_seconds} & \code{time.perf\_counter()}; tiempo de ejecución \\
\code{action} & Acción elegida (primera en orden con score estrictamente mayor) \\
\code{value} & Valor Minimax de la raíz \\
\code{root\_values} & Tupla \code{(acción, valor)} por acción raíz; en $\alpha$-$\beta$ los posteriores a la mejor pueden ser \textbf{cotas} \\
\code{considered} & (si \code{record=True}) Tupla \code{(GameState, Action)} de cada acción considerada (check 6) \\
\bottomrule
\end{tabular}
\end{table}
\label{tab:metricas-adversarial}


\section{RF-31: Misma Decisión con Mismo Orden}
\label{sec:rf31}

En la raíz, empate en score $\to$ se mantiene la \textbf{primera acción en el orden dado}.
Tanto Minimax como $\alpha$-$\beta$ iteran las acciones raíz en el mismo orden
(\code{default} o \code{reversed}) y actualizan \code{chosen} solo si \code{score > best\_value}
(estrictamente mayor). Por tanto:

\begin{itemize}
\item Misma profundidad, misma utilidad, mismo orden $\to$ \textbf{misma decisión}
\item Valores de acciones posteriores a la mejor en $\alpha$-$\beta$ pueden ser \textbf{cotas}
  (no valores exactos), porque el corte impide explorar completamente
\end{itemize}

\section{RF-32: Orden de Acciones y Poda}
\label{sec:rf32}

Orden de acciones configurable: \code{default} (arriba/abajo/izq/der por unidad) o
\code{reversed}. Cambiar el orden:

\begin{itemize}
\item \textbf{No cambia la decisión} (RF-31) si la utilidad y profundidad son iguales
\item \textbf{Sí cambia \code{pruned} y \code{generated}} (RF-32)
\item En caso 9 (profundidad 3): \code{default} genera 142, \code{reversed} genera 191
\item En caso 9 (profundidad 4): orden \code{reversed} cambia el tie-break entre valores
  iguales (0), produciendo acción diferente con mismo valor
\end{itemize}

\section{Check 6 del Verificador (RF-42)}
\label{sec:check6}

\code{verify\_adversarial(board, state, result)} re-ejecuta con \code{record=True} y valida:

\begin{enumerate}
\item Cada acción en \code{result.considered} es legal según \code{legal\_actions}
\item La acción elegida (\code{result.action}) es legal
\item No se generan estados inválidos durante la exploración
\end{enumerate}

Esto completa los 6 checks de RF-42 (5 para búsquedas + 1 para adversarial).

\section{Casos 7--10: Resultados}
\label{sec:casos7-10}

Escenario: \code{caso\_07\_decision\_adversarial.json} (5$\times$7, coste 1, 2 unidades/bando,
recurso (0,1), A mueve primero, pesos 10/3/1/0.5).

\begin{table}[H]
\centering
\caption{Casos 7--10: Resultados Minimax vs Alfa-Beta}
\begin{tabular}{lccccccc}
\toprule
Alg & Prof & Orden & Acción & Valor & Gen & Eval & Podados \\
\midrule
MM & 1 & def & A2$\to$(0,1) & 3 & 7 & 7 & 0 \\
MM & 2 & def & A2$\to$(0,1) & 0 & 71 & 64 & 0 \\
MM & 3 & def & A1$\to$(1,1) & 3 & 538 & 467 & 0 \\
MM & 4 & def & A1$\to$(0,0) & 0 & 4,826 & 4,296 & 0 \\
AB & 1 & def & A2$\to$(0,1) & 3 & 7 & 7 & 0 \\
AB & 2 & def & A2$\to$(0,1) & 0 & 31 & 24 & 40 \\
AB & 3 & def & A1$\to$(1,1) & 3 & 142 & 110 & 106 \\
AB & 4 & def & A1$\to$(0,0) & 0 & 431 & 299 & 757 \\
AB & 3 & rev & A1$\to$(1,1) & 3 & 191 & 148 & 145 \\
\bottomrule
\end{tabular}
\end{table}
\label{tab:casos7-10}


\begin{itemize}
\item \textbf{Caso 7:} La mejor decisión no es ir directo al recurso (profundidad 1 y 2 agarran;
  profundidad 3 protege moviendo A1$\to$(1,1) para bloquear a B1).
\item \textbf{Caso 8:} En profundidad 4, $\alpha$-$\beta$ genera \textbf{11$\times$ menos nodos}
  que Minimax (431 vs 4,826).
\item \textbf{Caso 9:} Mismo valor (3), misma decisión, pero orden \code{reversed} poda más
  (191 vs 142 generados en d=3). En d=4, orden \code{reversed} cambia tie-break (acción
  diferente, mismo valor 0).
\item \textbf{Caso 10:} Profundidades impares terminan en turno de A (valor 3, guarda);
  pares terminan en turno de B (valor 0, aplana). Efecto horizonte documentado en \S\ref{sec:limitaciones-conocidas}.
\end{itemize}